\chapter{Wenn das Modell versagt}\label{chap:modell}

\einstieg{Ein Modell sagt: Ein solcher Tag kommt alle 13'800 Jahre. Dann passiert er zweimal in zehn Jahren. Wer haftet?}

\arbeitsblock{Notebook \texttt{07\_Modell.ipynb}, 35 Minuten. Danach 25 Minuten Projektarbeit.}

\section{Dicke Ränder}
Bisher haben wir meist normalverteilte Renditen angenommen. Reale Kursbewegungen zeigen aber häufiger extreme Tage, als die Normalverteilung erwartet. Wir vergleichen zwei Modelle mit gleicher Tagesvolatilität von 1~\% über 2520 Handelstage, also rund zehn Jahre. Das zweite Modell verwendet eine t-Verteilung mit 3 Freiheitsgraden.

\begin{table}[H]
	\centering
	\small
	\caption{Durchschnitt aus 2000 simulierten Zehnjahres-Perioden}
	\label{tab:extrem}
	\begin{tabular}{@{}lrrrr@{}}
		\midrule
		\textbf{Modell} & \textbf{Tage $\leq -3$ \%} & \textbf{Tage $\leq -5$ \%} & \textbf{\acrshort{VaR} 99 \%} & \textbf{Schlechtester Tag} \\
		\midrule
		Normalverteilung & 3.5 & 0.0 & $-2.3$ \% & $-3.4$ \% \\
		t-Verteilung & 17.2 & 4.0 & $-2.6$ \% & $-9.0$ \% \\
		\midrule
	\end{tabular}
\end{table}

Der \gls{VaR} 99~\% gibt den Tagesverlust an, der nur an einem von 100 Tagen überschritten wird. Beide Modelle liegen dabei nahe beieinander, obwohl sich die Extremtage deutlich unterscheiden. Der \acrshort{VaR} sagt nichts darüber, wie schlimm die übrigen Tage ausfallen. Beim schlechtesten Tag ist der Median angegeben.

\begin{myexercise}[Fünf Standardabweichungen]
	Ein Verlust von 5~\% entspricht hier fünf Standardabweichungen. Die Normalverteilung ordnet dem eine Wahrscheinlichkeit von rund $2.9\cdot 10^{-7}$ zu. Wie oft käme ein solcher Tag bei 252 Handelstagen pro Jahr vor?
	\begin{myanswer}
		Etwa einmal in $1/(2.9\cdot 10^{-7}\cdot 252)$, also rund 13'800 Jahren. Wenn reale Märkte solche Tage innerhalb weniger Jahre zeigen, liegt das Problem beim Modell.
	\end{myanswer}
\end{myexercise}

\begin{mychallenge}[Schlimmer als der Value at Risk]
	Der Expected Shortfall 99~\% ist der Durchschnitt aller Tage, die schlechter sind als der \acrshort{VaR}. Berechnen Sie beide Zahlen für beide Modelle.
	\hilfe{\texttt{v = np.percentile(r, 1)} liefert den \acrshort{VaR}, \texttt{r[r <= v].mean()} den Expected Shortfall.}
	\begin{myanswer}
		Der \acrshort{VaR} liegt bei $-2.3$~\% und $-2.6$~\%, der Expected Shortfall bei $-2.7$~\% und $-4.0$~\%. Der Expected Shortfall trennt die Modelle deutlich besser, weil er in den Rand hineinschaut.
	\end{myanswer}
\end{mychallenge}

\begin{mychallenge}[Schulden als Verstärker]
	Ein Fonds hat eine Tagesvolatilität von 0.5~\%, eine Tagesdrift von 0.04~\% und arbeitet mit Hebel $h$. Wie oft verliert er innerhalb eines Jahres 80~\% seines Eigenkapitals?
	\hilfe{Die Tagesrendite des Eigenkapitals ist $h \cdot r$. Rechnen Sie mit Logarithmen, denn die lassen sich mit \texttt{np.cumsum} aufsummieren. Im Notebook steht der passende Ansatz. Vergleichen Sie die Summe mit \texttt{np.log(0.2)}.}
	\begin{myanswer}
		\begin{center}
			\small
			\begin{tabular}{@{}lrrr@{}}
				\midrule
				\textbf{Modell} & \textbf{Hebel 5} & \textbf{Hebel 10} & \textbf{Hebel 25} \\
				\midrule
				Normalverteilung & 0.0 \% & 0.5 \% & 31 \% \\
				t-Verteilung & 0.2 \% & 1.9 \% & 34 \% \\
				\midrule
			\end{tabular}
		\end{center}
		Bei Hebel 10 liegt das Risiko mit dicken Rändern rund viermal höher. Bei Hebel 25 reicht schon die normale Schwankung. Der Hebel wirkt stärker als die Wahl der Verteilung.
	\end{myanswer}
\end{mychallenge}

\section{Fallbeispiele}
Die SchülerInnen wählen in Gruppen einen Fall und klären drei Fragen: Was hat das Modell angenommen, was ist tatsächlich passiert, und wer trug den Schaden?

\begin{itemize}
	\item \textbf{Frankenschock, Januar 2015:} Die Schweizerische Nationalbank hebt den Mindestkurs zum Euro auf. Ein Kurs, der lange fast fest war, springt innert kurzer Zeit. Was misst eine Volatilität, die vor dem Ereignis geschätzt wurde?
	\item \textbf{Long-Term Capital Management, 1998:} Ein Hedgefonds mit renommierten Finanzökonomen arbeitet mit hohen Schulden und gerät nach der Russlandkrise in Schieflage. Banken organisieren die Rettung. Weshalb verstärken Schulden Modellfehler?
	\item \textbf{Finanzkrise 2008:} Hypothekenkredite galten durch Bündelung als gut gestreut. Dann fielen viele gleichzeitig aus. Das schliesst direkt an die Korrelation in Kapitel 5 an: Wie unabhängig waren die Risiken wirklich?
\end{itemize}

\begin{myexercise}[Modellkritik]
	Nennen Sie drei Annahmen aus Ihrer bisherigen Simulation, die in der Realität verletzt sein könnten, und testen Sie eine davon im Notebook.
	\begin{myanswer}
		Beispiele sind unabhängige Jahre, konstante Volatilität, normalverteilte Renditen, keine Kosten und keine Steuern, oder dass eine Anlage immer weiterbesteht.
	\end{myanswer}
\end{myexercise}
